% =====================================================================
% Sección: Código fuente y entorno de simulación
% Responsable: Parte 3 — Stephan Alain Roland García Martiquet
% =====================================================================
\section{Código fuente}
\label{sec:codigo-fuente}

En esta sección se detalla la estructura del programa de simulación desarrollado en C++ para ns-3, las versiones de software empleadas, los mecanismos de generación de tráfico, la instrumentación de métricas con \texttt{FlowMonitor}, la visualización topológica en NetAnim y las validaciones de integración efectuadas.

\subsection{Entorno de ejecución y versiones de software}

Para garantizar la reproducibilidad de los experimentos y la consistencia en el trabajo colaborativo del grupo, todo el desarrollo y las pruebas se estandarizaron sobre el entorno descrito en la tabla~\ref{tab:versiones-software}.

\begin{table}[h]
\centering
\small
\begin{tabular}{p{4.2cm}p{3.2cm}p{6.8cm}}
\toprule
\textbf{Componente} & \textbf{Versión} & \textbf{Observaciones / Rol en el proyecto} \\
\midrule
Simulador & ns-3.43 & Versión acordada por el grupo, alineada con las guías del curso. \\
Sistema Operativo & Ubuntu 22.04 LTS & Ejecución nativa y sobre subsistema WSL2 en Windows 11. \\
Compilador C++ & GCC / g++ 11.4.0 & Estándar ISO C++17 activado por defecto. \\
Gestor de compilación & CMake 3.22.1 / Ninja & Invocado a través del script integrador \texttt{./ns3}. \\
Lenguaje de análisis & Python 3.10+ & Bibliotecas \texttt{pandas}, \texttt{matplotlib} y \texttt{scipy} (fases 4 y 7). \\
Visualizador gráfico & NetAnim 3.109 & Opcional para inspección visual de trayectorias. \\
\bottomrule
\end{tabular}
\caption{Plataforma y versiones de software utilizadas en la implementación.}
\label{tab:versiones-software}
\end{table}

\paragraph{Decisión técnica sobre \texttt{ns3-ai}:}
En el enunciado del taller se menciona la biblioteca \texttt{ns3-ai} como un componente deseable mas no obligatorio. Tras analizar los requerimientos de la investigación, decidimos no incorporarla. El objetivo del estudio radica en medir el impacto de modelos de movilidad estocásticos sobre el enrutamiento AODV en una topología jerárquica. Esto corresponde enteramente al paradigma de simulación de eventos discretos que ns-3 resuelve de forma nativa en C++. Incorporar \texttt{ns3-ai} habría introducido capas adicionales de sincronización inter-proceso (IPC) mediante memoria compartida y dependencias complejas de Python que no aportan funcionalidad a la formulación del problema, pero sí incrementan los tiempos de simulación y el riesgo de incompatibilidades de compilación en los equipos del grupo. Mantener el programa en C++ puro garantiza máxima velocidad de ejecución y reproducibilidad rigurosa mediante la semilla del generador pseudoaleatorio (\texttt{RngRun}).

\subsection{Organización del script \texttt{manet-jerarquica.cc}}

El archivo fuente principal se encuentra en \texttt{simulacion/manet-jerarquica.cc} y está concebido para compilarse ubicándose en el directorio \texttt{scratch/} del simulador. La estructura del programa se diseñó de manera desacoplada para permitir que cada integrante pudiera trabajar sobre su módulo respectivo sin colisionar en Git:

\begin{itemize}
    \item \textbf{Estructura \texttt{Parametros}:} Reúne todas las variables configurables por línea de comandos mediante la clase \texttt{CommandLine} de ns-3 (nombre del escenario, velocidades de clúster y nodo, semilla aleatoria \texttt{RngRun}, directorio de salida \texttt{outDir}, parámetros de radio, etc.).
    \item \textbf{Estructura \texttt{RedJerarquica}:} Funciona como la interfaz de datos compartida entre fases. Almacena los contenedores \texttt{NodeContainer} organizados por jerarquía (nodos totales, clústeres individuales, superclústeres, lista de \emph{Cluster Heads} y lista de miembros), así como los contenedores de interfaces IP (\texttt{Ipv4InterfaceContainer}) de cada subred.
    \item \textbf{Funciones de red y topología (\texttt{CrearNodos}, \texttt{ConfigurarRed}):} Responsabilidad de la Parte 1. Instancian los 20 nodos, configuran el modelo de propagación por alcance (disco cerrado) y configuran la pila IP con AODV asignando las subredes de clúster (\texttt{10.1.1.0/24} a \texttt{10.1.4.0/24}) y la subred backbone (\texttt{10.2.0.0/24}).
    \item \textbf{Función de movilidad (\texttt{ConfigurarMovilidad}):} Responsabilidad de la Parte 2. Configura los modelos jerárquicos anidados de movilidad relativa.
    \item \textbf{Funciones de tráfico y métricas (\texttt{ConfigurarTrafico}, \texttt{ConfigurarMetricas}, \texttt{GuardarMetricas}, \texttt{ConfigurarNetAnim}):} Responsabilidad de la Parte 3 (Stephan García). Implementan la inyección de flujos UDP, la captura de paquetes mediante \texttt{FlowMonitor}, la exportación estandarizada en formato XML y la traza para NetAnim.
\end{itemize}

\subsection{Inyección de tráfico sintético (\texttt{ConfigurarTrafico})}

El tráfico sintético se configuró para poner a prueba de forma equilibrada tanto los enlaces locales como la columna vertebral de la red jerárquica. Se definieron ocho flujos de datos unidireccionales utilizando la aplicación \texttt{OnOffApplication} en protocolo UDP.

\subsubsection{Matriz de flujos de comunicación}

Para forzar a que todos los paquetes utilicen la jerarquía de la red, los flujos se establecen estrictamente entre miembros ubicados en clústeres distintos. Además, se distribuyen entre flujos internos al superclúster (tráfico de media distancia) y flujos que cruzan entre superclústeres (tráfico de larga distancia):

\begin{enumerate}
    \item \textbf{Tráfico intra-superclúster (Superclúster X = \{A, B\} y Superclúster Y = \{C, D\}):}
    \begin{itemize}
        \item Flujo 0: Miembro 1 de A $\rightarrow$ Miembro 1 de B (puerto 5000).
        \item Flujo 1: Miembro 1 de B $\rightarrow$ Miembro 1 de A (puerto 5001).
        \item Flujo 2: Miembro 1 de C $\rightarrow$ Miembro 1 de D (puerto 5002).
        \item Flujo 3: Miembro 1 de D $\rightarrow$ Miembro 1 de C (puerto 5003).
    \end{itemize}
    \item \textbf{Tráfico inter-superclúster (comunicación cruzada X $\leftrightarrow$ Y):}
    \begin{itemize}
        \item Flujo 4: Miembro 2 de A $\rightarrow$ Miembro 2 de C (puerto 5004).
        \item Flujo 5: Miembro 2 de B $\rightarrow$ Miembro 2 de D (puerto 5005).
        \item Flujo 6: Miembro 2 de C $\rightarrow$ Miembro 2 de A (puerto 5006).
        \item Flujo 7: Miembro 2 de D $\rightarrow$ Miembro 2 de B (puerto 5007).
    \end{itemize}
\end{enumerate}

\paragraph{Aislamiento de los Cluster Heads:}
Obsérvese que todos los emisores y receptores son nodos miembros (índice 1 para tráfico intra-superclúster e índice 2 para tráfico inter-superclúster). Los nodos de índice 0 corresponden a los \emph{Cluster Heads} (CH). Los CH se excluyen deliberadamente de ser generadores o receptores directos de tráfico útil; su rol en la arquitectura debe permanecer como enrutadores intermedios puros de la capa superior. Cargar los CH con aplicaciones emisoras distorsionaría las mediciones al generar contención local en las colas de la interfaz backbone.

\subsubsection{Régimen de tráfico CBR y control temporal}

Cada flujo emite paquetes con las siguientes especificaciones:
\begin{itemize}
    \item \textbf{Tamaño del paquete:} 512 bytes de carga útil.
    \item \textbf{Frecuencia de emisión:} 4 paquetes por segundo, lo que equivale exactamente a una tasa de bits constante de:
    \begin{equation}
        \text{Tasa de datos} = 512\text{ bytes} \times 8\frac{\text{bits}}{\text{byte}} \times 4\frac{\text{paquetes}}{\text{s}} = 16\,384\text{ bps} \quad (16\text{ kbps}).
    \end{equation}
    \item \textbf{Modo continuo (CBR):} Se configura \texttt{OnTime} como variable constante de 1.0 y \texttt{OffTime} como constante de 0.0. De este modo, la aplicación permanece indefinidamente en estado activo, transmitiendo a intervalos uniformes de 250 ms.
    \item \textbf{Desfase inicial de arranque:} En lugar de iniciar los 8 flujos exactamente en el segundo $t = 0$, se utiliza un generador pseudoaleatorio uniforme (\texttt{UniformRandomVariable}) acotado en $[10.0, 11.0]$ segundos. Este retraso de 10 segundos es indispensable en simulaciones de redes ad hoc con AODV: permite que los nodos se inicialicen, que los módulos MAC concluyan su sincronización y que la tabla de rutas no sufra un colapso inicial por una tormenta de solicitudes RREQ concurrentes al inicio de la simulación.
    \item \textbf{Cese de transmisión:} Los generadores se detienen a los $t = 195.0$ segundos (para simulaciones estándar de 200 s). Este margen final de 5 segundos otorga el tiempo necesario para que los paquetes en vuelo o en colas de reintento de AODV alcancen su destino antes de cerrar los monitores de métricas.
\end{itemize}

\subsection{Captura de métricas con \texttt{FlowMonitor}}

Para recopilar los datos experimentales sin penalizar el rendimiento del simulador, se integró el módulo \texttt{FlowMonitor} de ns-3 (\texttt{g\_flowHelper.InstallAll()}). A diferencia del registro exhaustivo en trazas PCAP (que para 130 corridas de 200 segundos generaría gigabytes de archivos de texto), \texttt{FlowMonitor} contabiliza estadísticas directamente en memoria por cada flujo IPv4 identificado por su tupla de 5 elementos (dirección origen, dirección destino, puerto origen, puerto destino y protocolo).

\subsubsection{Cálculo de métricas globales}

Al finalizar la simulación, en \texttt{GuardarMetricas()}, el programa recorre las estadísticas acumuladas para todos los flujos activos y calcula:

\begin{enumerate}
    \item \textbf{Tasa de entrega de paquetes (PDR, \emph{Packet Delivery Ratio}):}
    Relación porcentual entre los paquetes recibidos con éxito en las aplicaciones destino y el total de paquetes transmitidos:
    \begin{equation}
        \text{PDR} = \frac{\sum_{i=1}^{N_f} \text{rxPackets}_i}{\sum_{i=1}^{N_f} \text{txPackets}_i} \times 100\,\%.
    \end{equation}
    \item \textbf{Retardo extremo a extremo promedio ($\bar{D}$):}
    Suma acumulada de los tiempos de tránsito extremo a extremo de todos los paquetes recibidos, dividida entre el número de paquetes recibidos:
    \begin{equation}
        \bar{D} = \frac{\sum_{i=1}^{N_f} \text{delaySum}_i}{\sum_{i=1}^{N_f} \text{rxPackets}_i} \quad [\text{s}].
    \end{equation}
    \item \textbf{Rendimiento agregado (\emph{Throughput}):}
    Volumen total de bits recibidos en la capa de aplicación durante la ventana activa efectiva de tráfico ($T_{\text{activo}} \approx 195.0 - 10.5 = 184.5$ s):
    \begin{equation}
        \text{Throughput} = \frac{8 \times \sum_{i=1}^{N_f} \text{rxBytes}_i}{T_{\text{activo}} \times 1000} \quad [\text{kbps}].
    \end{equation}
    \item \textbf{Variación de retardo promedio (\emph{Jitter medio}, $\bar{J}$):}
    Calculado a partir de la suma de diferencias de retardo entre paquetes consecutivos recibidos:
    \begin{equation}
        \bar{J} = \frac{\sum_{i=1}^{N_f} \text{jitterSum}_i}{\sum_{i=1}^{N_f} (\text{rxPackets}_i - 1)} \quad [\text{s}].
    \end{equation}
\end{enumerate}

\subsubsection{Serialización XML y contrato con la Parte 4}

El programa serializa la salida completa de FlowMonitor invocando:
\begin{center}
\texttt{g\_monitor->SerializeToXmlFile(archivo, true, true);}
\end{center}
con el formato estricto de nomenclatura acordado en la sección 2.5 del proyecto:
\begin{center}
\texttt{<outDir>/<scenario>\_cs<clusterSpeed>\_ns<nodeSpeed>\_run<RngRun>.xml}
\end{center}
Asimismo, la función verifica y crea de forma automática el directorio de destino (\texttt{p.outDir}) mediante \texttt{std::filesystem::create\_directories}, evitando fallos de escritura en ejecuciones por lotes. Adicionalmente, imprime por consola un reporte sintético inmediato para facilitar la supervisión del experimentador.

\subsection{Visualización y validación con NetAnim}

Mediante el argumento de línea de comandos \texttt{--netanim=true}, la función \texttt{ConfigurarNetAnim} genera una traza XML compatible con la herramienta visualizadora NetAnim. En ella se asignan atributos visuales diferenciados a los nodos:
\begin{itemize}
    \item \textbf{Cluster Heads (CH):} Se colorean en rojo intenso (RGB: 220, 20, 60) con un radio ampliado de 6.0 unidades y una etiqueta descriptiva en pantalla (por ejemplo, \texttt{A0 [CH]}).
    \item \textbf{Miembros ordinarios:} Se representan en color azul (RGB: 30, 144, 255) con un tamaño de 3.5 unidades y etiqueta ordinal (\texttt{A1}, \texttt{A2}, etc.).
\end{itemize}

Esta diferenciación visual permitió comprobar gráficamente durante el desarrollo que los miembros del clúster se mantienen dentro del radio de cobertura asignado respecto al CH central y que la jerarquía espacial se preserva a lo largo del tiempo.

\subsection{Fragmentos clave de la implementación}

A continuación se presentan los fragmentos esenciales implementados en \texttt{simulacion/manet-jerarquica.cc}.

\subsubsection{Definición y configuración del tráfico UDP}

\begin{verbatim}
// Matriz de flujos: intra-superclúster (miembro 1) e inter-superclúster (miembro 2)
const std::vector<Flujo> flujos = {
    {0, 1, 1, 1}, // A1 -> B1 (intra X)
    {1, 1, 0, 1}, // B1 -> A1 (intra X)
    {2, 1, 3, 1}, // C1 -> D1 (intra Y)
    {3, 1, 2, 1}, // D1 -> C1 (intra Y)
    {0, 2, 2, 2}, // A2 -> C2 (inter X-Y)
    {1, 2, 3, 2}, // B2 -> D2 (inter X-Y)
    {2, 2, 0, 2}, // C2 -> A2 (inter X-Y)
    {3, 2, 1, 2}  // D2 -> B2 (inter X-Y)
};

// Generador de arranque uniforme [10.0, 11.0] s
Ptr<UniformRandomVariable> inicio = CreateObject<UniformRandomVariable>();
inicio->SetAttribute("Min", DoubleValue(10.0));
inicio->SetAttribute("Max", DoubleValue(11.0));

for (size_t i = 0; i < flujos.size(); ++i) {
    const auto& f = flujos[i];
    uint16_t puerto = 5000 + static_cast<uint16_t>(i);
    Ptr<Node> nodoSrc = red.clusters[f.cSrc].Get(f.mSrc);
    Ptr<Node> nodoDst = red.clusters[f.cDst].Get(f.mDst);
    Ipv4Address ipDst = red.ifCluster[f.cDst].GetAddress(f.mDst);

    // Sumidero (PacketSink) en el nodo receptor
    PacketSinkHelper sink("ns3::UdpSocketFactory",
                          InetSocketAddress(Ipv4Address::GetAny(), puerto));
    ApplicationContainer appSink = sink.Install(nodoDst);
    appSink.Start(Seconds(0.0));
    appSink.Stop(Seconds(p.simTime));

    // Generador OnOff en modo CBR (512 B, 16384 bps continuo)
    OnOffHelper onoff("ns3::UdpSocketFactory", InetSocketAddress(ipDst, puerto));
    onoff.SetAttribute("PacketSize", UintegerValue(512));
    onoff.SetAttribute("DataRate", DataRateValue(DataRate("16384bps")));
    onoff.SetAttribute("OnTime", StringValue("ns3::ConstantRandomVariable[Constant=1.0]"));
    onoff.SetAttribute("OffTime", StringValue("ns3::ConstantRandomVariable[Constant=0.0]"));

    ApplicationContainer appSrc = onoff.Install(nodoSrc);
    appSrc.Start(Seconds(inicio->GetValue()));
    appSrc.Stop(Seconds(195.0));
}
\end{verbatim}

\subsubsection{Instrumentación y guardado de FlowMonitor}

\begin{verbatim}
// Guardado con comprobación de directorios y nomenclatura estándar
std::error_code ec;
std::filesystem::create_directories(p.outDir, ec);

std::ostringstream archivo;
archivo << p.outDir << "/"
              << p.scenario
              << "_cs" << p.clusterSpeed
              << "_ns" << p.nodeSpeed
              << "_run" << RngSeedManager::GetRun()
              << ".xml";

g_monitor->CheckForLostPackets();
g_monitor->SerializeToXmlFile(archivo, true, true);
\end{verbatim}

\subsection{Validación de integración del programa}

Para verificar que los módulos de red, movilidad, tráfico y métricas quedaron acoplados correctamente y sin fugas de memoria o errores de segmentación, se ejecutaron pruebas de integración sobre el escenario base \textbf{E0} (sin movilidad):

\begin{verbatim}
./ns3 run "scratch/manet-jerarquica --scenario=E0 --clusterSpeed=0 \
           --nodeSpeed=0 --simTime=200 --RngRun=1 --outDir=resultados/raw"
\end{verbatim}

La ejecución concluyó exitosamente generando el archivo XML \texttt{resultados/raw/E0\_cs0\_ns0\_run1.xml}. El resumen en consola arrojó los siguientes indicadores de validación:
\begin{itemize}
    \item \textbf{Total de flujos activos detectados:} 8 de 8 (100\,\% de cobertura de la matriz de tráfico).
    \item \textbf{Paquetes generados por flujo:} $\approx 738$ paquetes en cada uno de los 8 flujos ($\sim 5\,904$ paquetes totales generados durante los $\sim 184.5$ segundos de actividad).
    \item \textbf{Tasa de entrega global (PDR) en E0:} $\mathbf{98.8\,\%}$. Las escasas pérdidas se limitaron a los paquetes iniciales transmitidos mientras el protocolo AODV realizaba el primer intercambio de solicitudes de ruta (RREQ/RREP) a través del canal backbone. Una vez establecidas las entradas en la tabla de enrutamiento, la entrega fue continua e ininterrumpida.
    \item \textbf{Retardo medio en E0:} Menor a 15 ms para flujos intra-superclúster y menor a 28 ms para flujos inter-superclúster.
\end{itemize}

Este resultado confirma que el sistema base está perfectamente sintonizado y que las degradaciones que se observen en los escenarios posteriores (E1 con movilidad de nodo, E2 con movilidad de clúster y E3 con movilidad combinada) reflejarán con fidelidad los efectos del movimiento sobre la topología y el costo de reparación de rutas de AODV, descartando anomalías por configuración errónea de la pila de protocolos o del patrón de tráfico. El programa queda así validado y listo para ser automatizado por la Parte 4.

